\documentclass[11pt]{letter}
\usepackage[utf8]{inputenc}
\usepackage{geometry}
\usepackage{enumitem}
\geometry{a4paper, margin=1in}

\address{Shubashis Mete \\
Department of Computer Science and Engineering \\
SRM Institute of Science and Technology \\
Chennai, India \\
shubashismete@gmail.com}

\date{\today}

\begin{document}

\begin{letter}{To the Editor-in-Chief \\
Education and Information Technologies \\
Springer Nature}

\opening{Dear Editor-in-Chief,}

I am pleased to submit our original research manuscript entitled \textbf{``Adaptive Personalized Tutoring System using Hybrid Deep Reinforcement Learning and Semantic NLP''} for consideration for publication in \textit{Education and Information Technologies}.

Our research addresses the critical need for scalable, personalized digital learning environments that can adapt to the diverse cognitive and behavioral patterns of students. While many information technologies in education rely on static branching or simple heuristics, our work introduces a \textbf{Hybrid RL-NLP framework} that optimizes long-term pedagogical strategies while maintaining semantic understanding and prerequisite safety.

We believe our work aligns perfectly with the scope of \textit{Education and Information Technologies}, particularly in the areas of \textbf{Adaptive Learning Systems, AI in Education, and Educational Data Mining}. This manuscript will be of interest to the readers of EAIT because it provides a bridge between advanced computational architectures and practical pedagogical implementation, offering actionable insights for researchers and practitioners aiming to deploy safe, scalable, and personalized digital learning tools. The key highlights of our submission include:

\begin{itemize}[leftmargin=*, noitemsep]
    \item \textbf{Pedagogical Optimization:} A novel \textbf{Delta-Reward} function that aligns reinforcement learning agents with normalized knowledge gain, ensuring the tutor optimizes for mastery rather than short-term scores.
    \item \textbf{Natural Language Understanding:} Integration of a Transformer-based semantic module that allows the system to assess open-ended student responses with a $0.62$ correlation to human judgment, facilitating more natural human-computer interaction in learning.
    \item \textbf{System Reliability:} A safety-constrained layer that enforces domain-specific pedagogical rules, addressing the common concern of AI "hallucinations" or illogical sequencing in educational software.
    \item \textbf{Robust Empirical Evidence:} A large-scale controlled simulation study ($N=300$) calibrated on the OULAD dataset, showing a significant improvement in learning outcomes ($g=0.59$ vs $0.36$) with a large effect size of \textbf{Cohen's $d=1.50$}.
\end{itemize}

This study provides both technical contributions to the field of educational AI and practical insights for the design of the next generation of digital learning platforms. It bridges the gap between complex machine learning architectures and established educational theories such as the \textit{Zone of Proximal Development} and \textit{Mastery Learning}.

This manuscript is original work, has not been published elsewhere, and is not currently under consideration by any other journal. All authors have approved the submission. We have no conflicts of interest to disclose.

Thank you for your time and for considering our work for publication in \textit{Education and Information Technologies}.

\closing{Sincerely,}

\vspace{0.5cm}
Shubashis Mete \\
Corresponding Author \\
SRM Institute of Science and Technology

\end{letter}

\end{document}
